% !TEX root = ../main.tex
% !TEX program = lualatex
\documentclass[../main.tex]{subfiles}
\IfSubfilesClassLoaded{\externaldocument{\subfix{../build/main}}}{}
% =============================================================================
% File: 27C_Presentation_Techniques.tex
% =============================================================================
\begin{document}
\IfSubfilesClassLoaded{\chapter{EDO Study Guide}}{}
\section{Presentation Techniques}

\subsection{Summary}
\begin{description}
  \item[\textbf{Preparation.}] Effective briefs begin with audience analysis, understanding the deliverable, and clarifying the purpose (inform, inspire, motivate, persuade, or decide)~\autocite{edo-5-1-4-presentation-techniques-2025}.
  \item[\textbf{Structure.}] A clear opening, focused body, and decisive closing, with deliberate slide limits and time for questions, drives a coherent message~\autocite{edo-5-1-4-presentation-techniques-2025}.
  \item[\textbf{Delivery.}] Strong verbal, vocal, and visual execution, plus disciplined Q\&A handling, determines whether the audience trusts the message~\autocite{edo-5-1-4-presentation-techniques-2025}.
\end{description}

\subsection{Student Presentation Requirements}
\begin{description}
  \item[\textbf{Scenario.}] Teams collaborate with \acp{ipt} to brief a professional decision brief aligned to the Integrated Case Study objectives~\autocite{edo-5-1-4-presentation-techniques-2025}.
  \item[\textbf{Execution.}] Team briefs are limited to 20 minutes; each student briefs at least five slides and up to 10 minutes, with team grades for content and individual grades for presentation technique~\autocite{edo-5-1-4-presentation-techniques-2025}.
  \item[\textbf{Follow-up.}] Briefs are recorded for self-evaluation and debrief with the director of training~\autocite{edo-5-1-4-presentation-techniques-2025}.
\end{description}

\subsection{Preparation and Design}
\begin{description}
  \item[\textbf{Audience analysis.}] Identify the audience's background, expected questions, and desired outcomes to tailor depth and framing~\autocite{edo-5-1-4-presentation-techniques-2025}.
  \item[\textbf{Purpose.}] Define the purpose early and keep it concise; most Integrated Case Study briefs are decision briefs that seek buy-in~\autocite{edo-5-1-4-presentation-techniques-2025}.
  \item[\textbf{Structure.}] For briefs under 30 minutes, target roughly 15\% opening, 75\% body, 10\% closing, limit slide count, and plan three to five minutes per slide including questions~\autocite{edo-5-1-4-presentation-techniques-2025}.
  \item[\textbf{Opening.}] Gain attention, establish relevance and credibility, preview the topic, state purpose, and set two to three audience expectations tied to closing takeaways~\autocite{edo-5-1-4-presentation-techniques-2025}.
  \item[\textbf{Body.}] Present evidence for the purpose, use logical transitions, and structure content with chronological, problem/solution, or cause/effect patterns~\autocite{edo-5-1-4-presentation-techniques-2025}.
  \item[\textbf{Closing.}] Reinforce takeaways, make the final case for the recommended option, and ask for a decision without recapping the entire brief~\autocite{edo-5-1-4-presentation-techniques-2025}.
\end{description}

\subsection{Recommended Courses of Action}
\begin{description}
  \item[\textbf{Option framing.}] Present the recommended option first or last, describe all options objectively, and explain assumptions, risks, and impacts if funding is not received~\autocite{edo-5-1-4-presentation-techniques-2025}.
  \item[\textbf{Supporting material.}] Use charts, tables, and graphics sparingly, label everything, keep units consistent, and leverage backup slides for detail~\autocite{edo-5-1-4-presentation-techniques-2025}.
\end{description}

\subsection{Delivery Components}
\begin{description}
  \item[\textbf{Three components.}] Presentation effectiveness depends on verbal content, vocal delivery, and visual presence; all must be prepared together~\autocite{edo-5-1-4-presentation-techniques-2025}.
  \item[\textbf{Visual design.}] Use white space, minimal animation, fragments instead of full sentences, left-justified text, consistent colors with legends, readable fonts (about 16--24), and numbered slides; spell out acronyms and test visuals on the actual screen~\autocite{edo-5-1-4-presentation-techniques-2025}.
  \item[\textbf{Body language.}] Stand with balanced posture, move purposefully, make eye contact across the room, and avoid nervous gestures or overusing the pointer~\autocite{edo-5-1-4-presentation-techniques-2025}.
  \item[\textbf{Verbal and vocal habits.}] Speak clearly and concisely, avoid filler words and reading slides, use examples, command attention with authentic energy, and keep humor controlled~\autocite{edo-5-1-4-presentation-techniques-2025}.
\end{description}

\subsection{Q\&A and Tough Questions}
\begin{description}
  \item[\textbf{Q\&A discipline.}] Relax, pause, repeat questions if needed, start with the conclusion, then explain and stop when answered; ensure the audience understands and build in time for questions~\autocite{edo-5-1-4-presentation-techniques-2025}.
  \item[\textbf{Tough questions.}] Do not bluff; provide facts, identify missing data, offer qualified opinions, and keep the brief on track without letting non-principals derail the session~\autocite{edo-5-1-4-presentation-techniques-2025}.
\end{description}

\subsection{Virtual Presentation Tips}
\begin{description}
  \item[\textbf{Virtual delivery.}] Use the webcam, ensure bandwidth, minimize distractions, engage with polls or direct questions, and adjust visuals (larger fonts, lighting, professional background) to compensate for reduced presence~\autocite{edo-5-1-4-presentation-techniques-2025}.
\end{description}

\subsection{Common Failure Modes}
\begin{description}
  \item[\textbf{Death by PowerPoint.}] The failure is not the tool, but poor connection to the audience, unclear expectations, weak delivery, inaccurate content, or slides that are too dense or disorganized~\autocite{edo-5-1-4-presentation-techniques-2025}.
\end{description}

\subsection{Checklist and Written Communication}
\begin{itemize}
  \item Ensure smooth slide transitions, clear expectations and takeaways, accurate numbers, labeled charts, consistent formatting, spelled-out acronyms, tested equipment, and time for questions~\autocite{edo-5-1-4-presentation-techniques-2025}.
  \item Written communication should be clear and concise, tailored to the receiver, fact-checked, professional, and explicit about response timelines~\autocite{edo-5-1-4-presentation-techniques-2025}.
\end{itemize}

\subsection{Quick Review}
\begin{itemize}
  \item Analyze the audience, understand the deliverable, and define the purpose before building slides~\autocite{edo-5-1-4-presentation-techniques-2025}.
  \item Opening sets expectations, the body supports them, and the closing drives the decision or call to action~\autocite{edo-5-1-4-presentation-techniques-2025}.
  \item Verbal, vocal, and visual execution must align to earn trust and drive outcomes~\autocite{edo-5-1-4-presentation-techniques-2025}.
\end{itemize}

\IfSubfilesClassLoaded{
  \RenewDocumentCommand{\entryname}{}{\textbf{\color{Modern} Acronym}}
  \RenewDocumentCommand{\descriptionname}{}{\textbf{\color{Modern} Definition}}
  \printnoidxglossary[
    type=\acronymtype,
    title=Acronyms,
    style=long-booktabs
  ]}{}
\end{document}
